\section{Giới thiệu đề tài}

\subsection{Đặt vấn đề và mục tiêu thiết kế}

Phép nhân ma trận là phép toán cơ bản trong nhiều bài toán tính toán số.
Khi thực hiện trên FPGA, bộ nhân cần cho kết quả chính xác, có thời gian
xử lý phù hợp và sử dụng tài nguyên phần cứng hợp lý.

Mỗi phần tử $C_{i,j}$ được tính bằng tích vô hướng giữa hàng $i$ của
ma trận $A$ và cột $j$ của ma trận $B$. Với hai ma trận $8\times8$,
bộ nhân cần thực hiện $64\times8=512$ phép nhân để tạo đủ 64 phần tử
của $C$. Mỗi phần tử của $A$ được dùng cho tám cột kết quả, còn mỗi
phần tử của $B$ được dùng cho tám hàng kết quả. Đặc điểm này cho phép
tái sử dụng toán hạng đã đọc từ bộ nhớ.

Đề tài hướng đến các mục tiêu sau:

\begin{itemize}
    \item \textbf{Thực hiện chính xác phép nhân ma trận.}
    Thiết kế bộ nhân $C=A\cdot B$ trên FPGA với hai ma trận đầu vào
    $8\times8$, phần tử có dấu 16 bit.

    \item \textbf{Tối ưu phép nhân ma trận.}
    Tối ưu nhân ma trận với data reuse và kiểm thử so sánh với những
    phương pháp tuần tự khác.

    \item \textbf{Đánh giá chi phí phần cứng.}
    So sánh số phần tử logic, tài nguyên sử dụng trên bộ nhân FPGA 
    và công suất tiêu thụ trên bộ nhân.
\end{itemize}

\subsection{Phạm vi triển khai và tiêu chí đánh giá}

\begin{itemize}
    \item \textbf{Phạm vi triển khai:} Thiết kế bộ nhân hai ma trận
    $8\times8$ trên FPGA Cyclone IV E EP4CE6E22C8.

    \item \textbf{Độ chính xác:} Bộ nhân phải cho kết quả chính xác
    trên toàn bộ miền giá trị đầu vào 16 bit có dấu và đầu ra 16 bit
    có dấu và đối chiếu so sánh với nhân ma trận thủ công.

    \item \textbf{Đánh giá:} So sánh số chu kỳ xử lý, thông lượng,
    số LE, số bộ nhân chuyên dụng, số khối nhớ M9K.
\end{itemize}
